\documentclass[11pt]{article}

\usepackage[margin=1in]{geometry}
\usepackage[T1]{fontenc}
\usepackage{lmodern}
\usepackage{amsmath}
\usepackage{amssymb}
\usepackage{booktabs}
\usepackage{multirow}
\usepackage{graphicx}
\usepackage{rotating}
\usepackage{lscape}
\usepackage[colorlinks=true,linkcolor=blue]{hyperref}
\usepackage{cleveref}

\title{Rotated Floats and Landscape Layout: A Stress Test for Page Orientation}
\author{Test Corpus Author}
\date{\today}

\begin{document}

\maketitle

\begin{abstract}
This note exercises page-orientation features supplied by the \texttt{rotating}
and \texttt{lscape} packages. It contains a sideways figure float, a full
landscape block enclosing a wide multi-column table, and a vertically rotated
text span inside an otherwise upright table cell. Each rotated object carries a
caption or heading and a label, and the prose cross-references every float so
that a document-conversion pipeline must account both for the rotation and for
the numbering of the rotated objects.
\end{abstract}

\section{Motivation}
\label{sec:motivation}

Wide tabular material and large schematic figures frequently exceed the printable
width of a portrait page. The conventional remedy rotates the offending object
through ninety degrees, either as a floating \texttt{sidewaysfigure} or
\texttt{sidewaystable}, or by switching the entire page to landscape orientation
with a \texttt{landscape} environment. A third, more local device rotates a short
text span in place, which is useful for narrow column headings. We demonstrate
all three below and refer to each from the surrounding discussion.

The sideways schematic in \cref{fig:rotated} occupies the full text height when
rotated, which is why such figures are reserved for diagrams that would otherwise
be illegibly compressed. Its placeholder body is drawn with framing rules so that
no external image file is required.

\begin{sidewaysfigure}
  \centering
  \framebox[0.85\textwidth]{%
    \rule{0pt}{4cm}%
    \begin{minipage}{0.8\textwidth}
      \centering
      \rule{0.9\linewidth}{0.4pt}\\[1ex]
      \textbf{Schematic placeholder: data-processing pipeline}\\[1ex]
      \rule{0.9\linewidth}{0.4pt}\\[2ex]
      Ingest $\longrightarrow$ Clean $\longrightarrow$ Transform
      $\longrightarrow$ Model $\longrightarrow$ Report\\[2ex]
      \rule{0.9\linewidth}{0.4pt}
    \end{minipage}%
  }
  \caption{A sideways figure float rotated ninety degrees. The framed body
    stands in for a wide schematic and uses only rules and text, so the fixture
    needs no external graphics file.}
  \label{fig:rotated}
\end{sidewaysfigure}

\section{A wide table on a landscape page}
\label{sec:landscape}

When a table has many columns the cleanest solution is to devote an entire
landscape page to it. \Cref{tab:wide} reports a synthetic panel of summary
statistics across eight measurement columns; on a portrait page these columns
would collide with the right margin, whereas in the rotated landscape block they
sit comfortably within the rotated text width. The table uses \texttt{booktabs}
rules throughout.

\begin{landscape}
  \begin{table}
    \centering
    \caption{A wide eight-column panel of synthetic summary statistics, set on a
      landscape page so that all columns fit within the rotated text width.}
    \label{tab:wide}
    \begin{tabular}{l rrrr rrrr}
      \toprule
      & \multicolumn{4}{c}{Pre-treatment} & \multicolumn{4}{c}{Post-treatment} \\
      \cmidrule(lr){2-5} \cmidrule(lr){6-9}
      Region & Mean & S.D. & Min & Max & Mean & S.D. & Min & Max \\
      \midrule
      North      & 12.4 & 3.1 & 5.2 & 21.8 & 14.9 & 3.4 & 6.0 & 25.1 \\
      South      & 10.8 & 2.7 & 4.1 & 18.6 & 13.2 & 2.9 & 5.3 & 22.0 \\
      East       & 15.6 & 4.0 & 6.8 & 27.3 & 17.1 & 4.2 & 7.4 & 29.8 \\
      West       & 11.1 & 2.5 & 4.9 & 17.4 & 12.7 & 2.6 & 5.1 & 19.9 \\
      Central    & 13.7 & 3.3 & 5.7 & 23.0 & 15.8 & 3.6 & 6.5 & 26.4 \\
      \midrule
      All        & 12.7 & 3.1 & 4.1 & 27.3 & 14.7 & 3.3 & 5.1 & 29.8 \\
      \bottomrule
    \end{tabular}
  \end{table}
\end{landscape}

\section{Rotated text within an upright cell}
\label{sec:rotbox}

A finer-grained form of rotation turns a single span of text through ninety
degrees without rotating the surrounding page or float. \Cref{tab:rotbox} places
a vertically rotated row label in the left column of an otherwise ordinary
upright table, a layout often used when a long category name must head a narrow
column. The rotation is produced by \verb|\rotatebox{90}{...}| applied to the
cell contents.

\begin{table}[ht]
  \centering
  \caption{An upright table whose left-hand row label is rotated ninety degrees
    in place with \texttt{rotatebox}, while the remaining cells are set
    horizontally.}
  \label{tab:rotbox}
  \begin{tabular}{c l r}
    \toprule
    & Indicator & Value \\
    \midrule
    \multirow{3}{*}{\rotatebox{90}{\textbf{Quarter 1}}}
      & Revenue (\$m)      & 48.2 \\
      & Operating margin   & 0.19 \\
      & Headcount          & 312  \\
    \bottomrule
  \end{tabular}
\end{table}

\section{Discussion}
\label{sec:discussion}

The three devices serve complementary purposes. A \texttt{sidewaysfigure} such as
\cref{fig:rotated} rotates a floating object while leaving the page in portrait
orientation, which keeps the running text upright on the same physical page. A
\texttt{landscape} block, used for \cref{tab:wide}, rotates the entire page and is
therefore appropriate when the wide object should stand alone. An in-cell
\verb|\rotatebox|, illustrated by the row label of \cref{tab:rotbox}, rotates only
a short span and integrates with surrounding horizontal text. A conversion
pipeline must preserve the caption text, the label-based numbering, and the
intended orientation of each of \cref{fig:rotated}, \cref{tab:wide}, and
\cref{tab:rotbox}.

\end{document}
